\subsection{A finite translation-centralizer comparison at three and nine}

\begin{result}{A.4.8}{Proved}{The eighty-one translation-centralizer residues reduce to the geometric twenty-seven}
Let $E=F(a,b)$, let
$H=(\mathbb Z/3)^2$, and use the literal covering homomorphism
\[
 E\twoheadrightarrow H,\qquad a\longmapsto(1,0),
 \quad b\longmapsto(0,1),\qquad \Pi=\ker(E\to H).
                                                               \tag{TFC.1}
\]
It is the once-punctured-torus covering with nine labelled punctures;
$\Pi$ is free of rank ten.  Use the surface basis specified below.
Let $\mathcal P_m$ be the actual image of the pure nine-point genus-one
mapping class group on $H_1(\Pi,\mathbb Z/m)$, for $m=3,9$.
The two deck translations have the integral matrices $h_a,h_b$ in
(TFC.4), and
\[
 \mathcal P_m=\left\{
 \begin{pmatrix}C&0\\B&I_8\end{pmatrix}:
 C\in\operatorname{SL}_2(\mathbb Z/m),\quad
 B\in M_{8,2}(\mathbb Z/m)\right\}.                       \tag{TFC.2}
\]
Their simultaneous centralizer in $\mathcal P_3$ has exactly $81$
elements.  Precisely $27$ of these are reductions of elements in
$C_{\mathcal P_9}(H)$.  Those $27$ are exactly the homological image
modulo three of the actual geometric normalizer subgroup $N_0$
whose induced permutation of the nine labels is trivial.
The larger centralizer $C_{\mathcal P_9}(H)$ has order $3^7$.

This is a finite centralizer-image statement.  It does not assert
that all finite conjugates of $H$ in a larger mapping-class quotient
have already been geometrically straightened.
\end{result}

\begin{proof}
We first specify the maps and basis in (TFC.1).  Write $c=[a,b]$.
Take $A=b^3$, $B=a^3$ and eight puncture generators
$C_i=p_i c p_i^{-1}$ with the following literal words and labels:
\[
\begin{array}{c|c|c}
 i&p_i&q(p_i)\\ \hline
1&a^3b^3a^{-3}b^{-1}a&(1,2)\\
2&a^3b^3a^{-2}b^{-1}a^{-1}&(0,2)\\
3&a^2&(2,0)\\
4&a&(1,0)\\
5&1&(0,0)\\
6&ba^2&(2,1)\\
7&ba&(1,1)\\
8&b&(0,1).
\end{array}                                                   \tag{TFC.3}
\]
The ninth peripheral word, with label $(2,2)$, is
$b^2a^2c a^{-2}b^{-2}=([A,B]C_1\cdots C_8)^{-1}$.
These are actual free-basis substitutions.  One direct verification
uses the $3$-by-$3$ Cayley graph of $H$, with section
$t_{(i,j)}=b^ja^i$.  Take as spanning-tree edges the six horizontal
edges starting at $i=0,1$ and the two vertical edges starting at
$(0,0),(0,1)$.  The ten complementary edges give the Schreier basis.
Substituting (TFC.3), or successively eliminating the eight square
faces other than $(2,2)$, gives mutually inverse substitutions
between this Schreier basis and $A,B,C_1,\ldots,C_8$.
The executed program retains both substitutions and verifies them
as freely reduced words.  Thus no homology basis has been substituted
for a free group basis.  The downstairs closed period lattice is
$\Lambda=\mathbb Z[a]\oplus\mathbb Z[b]$, and the upstairs closed
period lattice is the actual sublattice $\Lambda'=3\Lambda$.
Ordinary forgetting upstairs maps $(A,B)$ to $(3b,3a)$ in
$\Lambda'$.  In its primitive ordered basis $(3a,3b)$ this map has
matrix $J=\left(\begin{smallmatrix}0&1\\1&0\end{smallmatrix}\right)$.
The covering map $\Lambda'\hookrightarrow\Lambda$, written in
the respective bases $(3a,3b)$ and $(a,b)$, has matrix $3I$;
the composite therefore has matrix $3J$, not $J$.
These two maps are recorded separately in the certificate.
The determinant-one assertion in (TFC.2) concerns automorphisms
of the upstairs lattice itself.  Its intersection form changes
sign under the basis change $J$, but its preserving automorphisms
still have determinant one.  For a downstairs normalizer with
integral matrix $M$, restriction to $3\Lambda$ has the same
matrix $M$ in $(3a,3b)$, and hence matrix $J M J$ in $(A,B)$.
This is an integral restriction of lattices before any reduction;
it does not divide a matrix modulo three by three.

Conjugation by $a,b$ on $\Pi$ gives the two deck actions.  In the
column basis $(A,B,C_1,\ldots,C_8)$ their integral matrices are
\[
 h_i=\begin{pmatrix}I_2&0\\t_i&P_i\end{pmatrix},
\quad
 t_a=\begin{pmatrix}0&0\\1&0\\0&0\\0&0\\1&0\\0&0\\0&0\\1&0\end{pmatrix},
\quad
 t_b=\begin{pmatrix}0&0\\0&0\\0&-1\\0&-1\\0&-1\\0&0\\0&0\\0&0\end{pmatrix},
                                                               \tag{TFC.4}
\]
where
\[
 P_a=\begin{pmatrix}
-1&1&0&0&0&0&0&0\\
-1&0&0&0&0&0&0&0\\
-1&0&0&1&0&0&0&0\\
-1&0&0&0&1&0&0&0\\
-1&0&1&0&0&0&0&0\\
-1&0&0&0&0&0&1&0\\
-1&0&0&0&0&0&0&1\\
-1&0&0&0&0&1&0&0
\end{pmatrix},\qquad
 P_b=\begin{pmatrix}
0&0&0&0&0&-1&1&0\\
0&0&0&0&0&-1&0&1\\
0&0&0&0&0&-1&0&0\\
1&0&0&0&0&-1&0&0\\
0&1&0&0&0&-1&0&0\\
0&0&1&0&0&-1&0&0\\
0&0&0&1&0&-1&0&0\\
0&0&0&0&1&-1&0&0
\end{pmatrix}.
\]
Direct multiplication verifies $h_a^3=h_b^3=I$ and $h_ah_b=h_bh_a$
over the integers.  Their permutation of all nine peripheral
classes is exactly addition of $(1,0)$ or $(0,1)$ to the labels in
(TFC.3).  The source program also checks the actual free-word
identity that the cube of each deck automorphism is the appropriate
inner automorphism by $a^3$ or $b^3$.

For (TFC.2), pure mapping classes fix each peripheral class and
induce an orientation-preserving torus homology map, so they lie
in the displayed group.  The converse is given by actual mapping
classes.  Two handle twists realize the elementary generators of
$\operatorname{SL}_2(\mathbb Z/m)$, and pushing the $i$th moving
puncture around the two handle loops realizes the two elementary
entries in the $i$th row of $B$.  For an explicit free-word check,
use the even kernel of $C_2(x_1)*\cdots*C_2(x_4)$ with basis
$u=x_2x_1,v=x_3x_1,w=x_4x_1$.  The surface basis
\[
 A=v,\quad B=u^{-1},\quad C=u^{-1}vuv^{-1}wv^{-1}u
\]
has boundary $x_1x_2x_3x_4=([A,B]C)^{-1}$.
The literal Hurwitz actions of the Ferrari pointpush basis
$\sigma_3\sigma_1^{-1}$ and
$\sigma_1\sigma_2\sigma_1\sigma_3^{-1}\sigma_2^{-1}\sigma_1^{-1}$
preserve this boundary and the conjugacy class of $C$; their
pointpush rows on $(A,B)$ are $(-1,0)$ and $(1,-1)$.
Their inverses and products therefore give both elementary rows
integrally.  Moving $C$ to each named puncture by the actual disk
braids gives the sixteen required generators.  This proves the
surjectivity assertion in (TFC.2), including its quotient maps.

A matrix $(C,B)$ in (TFC.2) centralizes the two translations exactly
when
\[
 (I-P_i)B=t_i(C-I),\qquad i=a,b.                         \tag{TFC.5}
\]
Over $\mathbb F_3$, solving these two displayed systems for all
$24$ values of $C\in\operatorname{SL}_2(\mathbb F_3)$ leaves only
$C=I$.  The linear map in the sixteen entries of $B$ has rank
$12$.  Its kernel has the following explicit parametrization:
\[
 u=\begin{pmatrix}2\\1\\1\\0\\2\\2\\1\\0\end{pmatrix},\qquad
 v=\begin{pmatrix}1\\2\\1\\2\\0\\2\\0\\1\end{pmatrix},\qquad
 B_0=(\alpha u+\gamma v,\ \beta u+\delta v),
 \quad(\alpha,\beta,\gamma,\delta)\in\mathbb F_3^4.
                                                               \tag{TFC.6}
\]
This gives all $81$ residues, without an omitted conjugator class.

Every lift of $C=I$ to $\operatorname{SL}_2(\mathbb Z/9)$ has
$C=I+3D$ with $\operatorname{tr}D=0$ over $\mathbb F_3$.
Writing $B=B_0+3B_1$, equation (TFC.5) becomes
\[
 (I-P_i)B_1=t_iD-\frac{(I-P_i)B_0}{3}\pmod3.             \tag{TFC.7}
\]
The division is performed on the literal integral matrices; it is
possible because (TFC.6) solves the equations modulo three.
The complete solution calculation shows that, before imposing the
trace condition, the uniquely permitted $D$ is
\[
 D(B_0)=\begin{pmatrix}
 2\alpha+\gamma&2\beta+\delta\\
 \alpha+\gamma&\beta+\delta
 \end{pmatrix}.                                        \tag{TFC.8}
\]
This can be verified on the four basis vectors in (TFC.6) and then
on all $81$ residues.  Changes of the chosen integral representative
of $B_0$ add an element in the image of the left side of (TFC.7),
so the resulting obstruction is linear over $\mathbb F_3$.
It follows that exactly the hyperplane
\[
 2\alpha+\beta+\gamma+\delta=0                          \tag{TFC.9}
\]
lifts to symplectic centralizers modulo nine.  It has $27$ elements.
For each, $D$ is unique and $B_1$ has the four-dimensional
homogeneous solution space (TFC.6).  This gives $27\cdot81=3^7$
elements in the larger centralizer.

Finally we identify this hyperplane with the actual geometric
normalizer image.  The normalizer fixing the origin label is
induced by the literal Nielsen automorphisms of $E$
\[
 U_E:(a,b)\mapsto(a,ba),\qquad
 L_E:(a,b)\mapsto(ab^{-1},b).
\]
They preserve $[a,b]$, preserve the characteristic covering kernel,
and permute its labels by their ordinary homology matrices.
The words $U_E^3,L_E^3,U_EL_E^3U_E^{-1}$ act trivially on all nine
labels.  Direct restriction to the basis (TFC.3) gives their four
coordinates in (TFC.6):
\[
 (2,0,2,0),\qquad(0,2,0,1),\qquad(0,0,1,2).
\]
They are independent and span (TFC.9).  All actual origin-fixing
normalizers with trivial label permutation lie over the principal
level-three subgroup of $\operatorname{SL}_2(\mathbb Z)$.
Its torsion-free Bass--Serre cover has rank three, so its image in
the elementary abelian group (TFC.6) has dimension at most three.
The displayed three elements attain that dimension.  Equivalently,
the executed finite Schreier computation of this image gives the
same three-dimensional span.  Thus the twenty-seven survivors are
precisely the stated geometric image.

The complete executed certificate is in
\path{translation_cover_basis.py},
\path{translation_mod9_centralizer.py}, and their JSON/NPZ
outputs.  All $81$ residues and their $27$ possible traceless closed
corrections are tested, giving $2187$ exact finite linear tests.
The auxiliary trace calculation verifies (TFC.8) on every residue.
\end{proof}

The remaining conjugacy issue is explicit.  In the actual image of
the pure group on the outer automorphisms of the class-two,
exponent-three quotient of $\Pi$, the outer IA kernel has dimension
$91$.  The ambient outer IA space has dimension $440$, and its
peripheral subspace has dimension $117$.  For the literal translation
subgroup, the kernel of
\[
 H^1(H,W_{91})\longrightarrow H^1(H,V_{117})
\]
has dimension six.  These are finite deformation directions, not
asserted genuine tower images.  The integral exponent-nine test
eliminates four of them and retains two at that stage.  The
centralizer result above consequently does not by itself establish
this conjugacy comparison.  A.4.9 below supplies the additional
characteristic degree-three peripheral lift and proves that its
actual class-two image removes all six of these directions.
